% from ega0-18-fr.tex
\oldpage[0\textsubscript{IV-1}]{67}

considéré comme un $(B/\mathfrak R)$-bimodule pour tout idéal ouvert
$\mathfrak R\subset\mathfrak R_0$. Soient alors
$\mathfrak R\subset\mathfrak R_0$ un idéal bilatère ouvert de $B$,
$\mathfrak J$ un idéal bilatère ouvert de $A$ tel que
$\rho(\mathfrak J)\subset\mathfrak R$, de sorte que $B/\mathfrak R$ peut être
considéré comme un $(A/\mathfrak J)$-anneau discret. La remarque précédente
prouve que le groupe
$\operatorname{Exan}_{A/\mathfrak J}(B/\mathfrak R,L)$ est défini~; en outre,
si $\mathfrak R'\subset\mathfrak R$, $\mathfrak J'\subset\mathfrak J$ sont
deux idéaux bilatères ouverts de $A$ et $B$ respectivement tels que
$\rho(\mathfrak J')\subset\mathfrak R'$, on a par (18.3.5.1) un homomorphisme
canonique
\begin{equation}
\label{0.18.5.1.1-fr}
\tag{18.5.1.1}
  \operatorname{Exan}_{A/\mathfrak J}(B/\mathfrak R,L)
  \longrightarrow
  \operatorname{Exan}_{A/\mathfrak J'}(B/\mathfrak R',L).
\end{equation}

L'ensemble des couples d'idéaux ouverts $(\mathfrak J,\mathfrak R)$ tels que
$\rho(\mathfrak J)\subset\mathfrak R\subset\mathfrak R_0$ est évidemment
ordonné filtrant à droite pour la relation «~$\mathfrak J\supset\mathfrak J'$
et $\mathfrak R\supset\mathfrak R'$~» et les applications (18.5.1.1)
définissent un \emph{système inductif} de groupes additifs ayant cet ensemble
pour ensemble d'indices. On pose, par abus de notation (car il ne s'agit plus
d'un groupe en correspondance biunivoque naturelle avec un ensemble
d'extensions)
\begin{equation}
\label{0.18.5.1.2-fr}
\tag{18.5.1.2}
  \operatorname{Exantop}_A(B,L)
  =\varinjlim\operatorname{Exan}_{A/\mathfrak J}(B/\mathfrak R,L).
\end{equation}

Dire que le second membre de (18.5.1.2) est \emph{nul} signifie donc que, pour
tout couple d'idéaux ouverts $\mathfrak J\subset A$, $\mathfrak R\subset B$
tels que $\rho(\mathfrak J)\subset\mathfrak R\subset\mathfrak R_0$ et toute
$(A/\mathfrak J)$-extension $E$ de $B/\mathfrak R$ par $L$, il existe deux
idéaux ouverts $\mathfrak J'\subset\mathfrak J$,
$\mathfrak R'\subset\mathfrak R$ tels que
$\rho(\mathfrak J')\subset\mathfrak R'$ et que l'image réciproque par
l'homomorphisme $B/\mathfrak R'\to B/\mathfrak R$ de $E$ soit \emph{triviale}.

Nous laissons au lecteur la définition analogue des limites inductives
$\operatorname{Exaltop}_A(B,L)$, $\operatorname{Exalcotop}_A(B,L)$ à partir de
$\operatorname{Exal}_{A/\mathfrak J}(B/\mathfrak R,L)$ et de
$\operatorname{Exalcom}_{A/\mathfrak J}(B/\mathfrak R,L)$ pour le cas où $A$
est commutatif et $B$ une $A$-algèbre (resp. une $A$-algèbre commutative)
topologique.
\end{env}

\begin{env}[18.5.2]
\label{0.18.5.2-fr}
Si l'on a un diagramme commutatif d'homomorphismes continus d'anneaux
\[
  \xymatrix@C=3pc@R=2pc{
    B'\ar[r] & B\\
    A'\ar[u]\ar[r] & A\ar[u]
  }
\]
on en déduit canoniquement deux homomorphismes de groupes additifs
\[
  \operatorname{Exantop}_A(B,L)
  \longrightarrow\operatorname{Exantop}_{A'}(B,L)
  \longrightarrow\operatorname{Exantop}_{A'}(B',L)
\]
par passage à la limite inductive à partir de (18.3.5.1).

En vertu de l'exactitude du foncteur $\varinjlim$ dans la catégorie des groupes
commutatifs, le noyau de l'homomorphisme
\[
  \operatorname{Exantop}_A(B,L)
  \longrightarrow\operatorname{Exantop}_{A'}(B,L)
\]
est la limite inductive des noyaux des homomorphismes
\[
  \operatorname{Exan}_{A/\mathfrak J}(B/\mathfrak R,L)
  \longrightarrow
  \operatorname{Exan}_{A'/\mathfrak J'}(B/\mathfrak R,L)
\]
où on a pris pour $\mathfrak J'$ l'image réciproque de $\mathfrak J$~; on note
ce noyau $\operatorname{Exantop}_{A/A'}(B,L)$. On définit de même
$\operatorname{Exaltop}_{A/A'}(B,L)$ et
$\operatorname{Exalcotop}_{A/A'}(B,L)$. Enfin, si on a un homomorphisme

